\documentclass[12pt, a4paper]{article}
%-----------------------
%- 	PACKAGES & SETTINGS
%-----------------------
\usepackage[utf8]{inputenc}
\usepackage[italian]{babel}
\usepackage{xcolor}
\usepackage{hyperref}
\hypersetup{
    colorlinks=true,
    filecolor=magenta,      
    urlcolor=blue,
    linkcolor=black
}
\urlstyle{same}
\usepackage{amsmath}
\usepackage{listings}
\lstset{language=c}
\usepackage{graphicx}
\graphicspath{ {images/} }
 
%-----------------------
%- 	TITLE
%-----------------------
\title{\textbf{Relazione Progetto Java PR 2}}
\author{\textbf{Venturi} Ludovico\\Docente: \href{http://pages.di.unipi.it/levi/}{Francesca Levi}}
\date{UNIPI, Novembre 2019}


%-----------------------
%- 	DOCUMENT
%-----------------------
\begin{document}
%- 	INTRO
\pagenumbering{roman} 
\maketitle
\tableofcontents
\vfill
\begin{figure}[h]
	\centering
	\includegraphics[scale=0.07]{javaLogo}
	\label{fig:0}
\end{figure}

\clearpage

%- 	START DOC
\pagenumbering{arabic} 

\section{Introduzione}
\begin{lstlisting}
	for i = 0 to n
		print "staminchia"
\end{lstlisting}
% \href{www.multiplayer.it}{MULTIPLAYER}
% figura \ref{fig:im2} \pageref{fig:im2}
 
%\begin{itemize}
%  \item The individual entries are indicated %with a black dot, a so-called bullet.
%  \item The text in the entries may be of any %length.
%\end{itemize}

%$\begin{enumerate}



%\begin{equation}
%E=mc^2
%\end{equation}
%\subsection{daie}
%\label{sec:daie1}
%$ E=mc^2 $
%$\Omega + 3 = 54$\\
%$\omega * 54$

%\ref{table:tab1} SEZIONE \ref{sec:daie1}

%h = here
%\begin{table}[h]
%\centering
%\begin{tabular}{|c|c|r|}
%	\hline
%		cell1 & gatto & gattini \\
%	\hline
%		cell1 & gatto & 12 \\
%	\hline
%		cell1 & gatto & gattini \\
%	\hline
%\end{tabular}
%\label{table:tab1}
%\end{table}


\end{document}
